\documentclass[11pt]{article}

\usepackage{tabularx}
\usepackage{array}
\usepackage{amssymb}   % needed for \blacksquare when you tick a box

% xcolor (loaded by macros_FP) and the later color package clash unless the
% options are declared up front; this keeps macros_FP.tex untouched.
\PassOptionsToPackage{usenames,dvipsnames}{color}
% start of paragraph with no indent
\parindent0pt

% insert packages in alphabetical order
\usepackage{algorithm, algpseudocode, amssymb, amsthm, bbm, csquotes, docmute, enumitem, float, geometry, graphicx, hyphenat, listings, mathtools, physics, tcolorbox, tikz, tocloft, xcolor, natbib}
\usepackage[english]{babel}

% load packages
\usepackage{geometry}
 \geometry{
 a4paper,
 total={170mm,257mm},
 left=20mm,
 top=20mm,
 }
\usepackage{amsfonts}
\usepackage[usenames, dvipsnames]{color}
%\usepackage{bbold}
%\usepackage{url}
\usepackage{upgreek}
\usepackage{placeins}
\usepackage{empheq}
\usepackage{tabularx}
\usepackage{booktabs}
\definecolor{red}{RGB}{255, 0, 0}
\usepackage{tabto}
\usepackage{tgtermes,tgheros} 
\usepackage{enumitem}% http://ctan.org/pkg/enumitem
\usepackage{titlesec}
\usepackage{hyperref}

\titleformat{\section}[block]
  {\fontsize{12}{15}\bfseries\sffamily\filcenter}
  {\thesection}
  {1em}
  {\MakeUppercase}
\titleformat{\subsection}[hang]
  {\fontsize{12}{15}\bfseries\sffamily}
  {\thesubsection}
  {1em}
  {}

\allowdisplaybreaks

% math symbols
\newcommand{\bA}{\mathbb{A}}
\newcommand{\bB}{\mathbb{B}}
\newcommand{\bC}{\mathbb{C}}
\newcommand{\bD}{\mathbb{D}}
\newcommand{\bE}{\mathbb{E}}
\newcommand{\bF}{\mathbb{F}}
\newcommand{\bG}{\mathbb{G}}
\newcommand{\bH}{\mathbb{H}}
\newcommand{\bI}{\mathbb{I}}
\newcommand{\bJ}{\mathbb{J}}
\newcommand{\bK}{\mathbb{K}}
\newcommand{\bL}{\mathbb{L}}
\newcommand{\bM}{\mathbb{M}}
\newcommand{\bN}{\mathbb{N}}
\newcommand{\bO}{\mathbb{O}}
\newcommand{\bP}{\mathbb{P}}
\newcommand{\bQ}{\mathbb{Q}}
\newcommand{\bR}{\mathbb{R}}
\newcommand{\bS}{\mathbb{S}}
\newcommand{\bT}{\mathbb{T}}
\newcommand{\bU}{\mathbb{U}}
\newcommand{\bV}{\mathbb{V}}
\newcommand{\bW}{\mathbb{W}}
\newcommand{\bX}{\mathbb{X}}
\newcommand{\bY}{\mathbb{Y}}
\newcommand{\bZ}{\mathbb{Z}}

\NewDocumentCommand{\codeword}{v}{%
\texttt{\textcolor{blue}{#1}}%
}


\newcommand{\cA}{\mathcal{A}}
\newcommand{\cB}{\mathcal{B}}
\newcommand{\cC}{\mathcal{C}}
\newcommand{\cD}{\mathcal{D}}
\newcommand{\cE}{\mathcal{E}}
\newcommand{\cF}{\mathcal{F}}
\newcommand{\cG}{\mathcal{G}}
\newcommand{\cH}{\mathcal{H}}
\newcommand{\cI}{\mathcal{I}}
\newcommand{\cJ}{\mathcal{J}}
\newcommand{\cK}{\mathcal{K}}
\newcommand{\cL}{\mathcal{L}}
\newcommand{\cM}{\mathcal{M}}
\newcommand{\cN}{\mathcal{N}}
\newcommand{\cO}{\mathcal{O}}
\newcommand{\cP}{\mathcal{P}}
\newcommand{\cQ}{\mathcal{Q}}
\newcommand{\cR}{\mathcal{R}}
\newcommand{\cS}{\mathcal{S}}
\newcommand{\cT}{\mathcal{T}}
\newcommand{\cU}{\mathcal{U}}
\newcommand{\cV}{\mathcal{V}}
\newcommand{\cW}{\mathcal{W}}
\newcommand{\cX}{\mathcal{X}}
\newcommand{\cY}{\mathcal{Y}}
\newcommand{\cZ}{\mathcal{Z}}
\newcommand{\vect}[1]{\mathbf{#1}}

\newcommand{\e}{\text{e}}

\newcommand{\ber}{\text{Ber}}
\newcommand{\bin}{\text{Bin}}
\newcommand{\poi}{\text{Poi}}
\newcommand{\geo}{\text{Geo}}
\newcommand{\Exp}{\text{Exp}}

\newcommand{\cov}{\text{Cov}}
\newcommand{\corr}{\rho}
\newcommand{\risk}{\text{R}}
\newcommand{\bias}{\text{Bias}}
\newcommand{\mse}{\text{MSE}}
\newcommand{\se}{\text{SE}}
\newcommand{\elbo}{\text{ELBO}}
\newcommand{\kl}{\text{KL}}

\newcommand{\acov}{\text{acov}}
\newcommand{\acorr}{\text{acorr}}
\newcommand{\xcov}{\text{xcov}}
\newcommand{\xcorr}{\text{xcorr}}
\newcommand{\pacorr}{\text{pacorr}}
\newcommand{\pacov}{\text{pacov}}

\newcommand{\relu}{\text{ReLU}}

\newcommand{\deriv}[2]{\frac{\partial #1}{\partial #2}}

\DeclareMathOperator*{\argmax}{arg\,max}
\DeclareMathOperator*{\argmin}{arg\,min}

% math environments
\theoremstyle{definition}
\newtheorem{definition}{Definition}[section]
\newtheorem{theorem}[definition]{Theorem}
\newtheorem{proposition}[definition]{Proposition}
\newtheorem{example}[definition]{Example}
\newtheorem{remark}[definition]{Remark}
\newtheorem{task}[definition]{Task}

\newenvironment{boxedtheorem}[1][]
{\begin{tcolorbox}[colback=white!95!gray, colframe=black, title=#1]}
{\end{tcolorbox}}

\newenvironment{boxeddefinition}[1][]
{\begin{tcolorbox}[colback=white, colframe=grey, title=#1]}
{\end{tcolorbox}}



\newcommand{\fillme}[1]{\textcolor{red}{[#1]}}

\begin{document}

\section*{AI Usage Disclosure}

\subsection*{Student information}

\begin{tabularx}{\textwidth}{>{\bfseries}l X}
Name: & Saliq Neyaz \\[0.6em]
Matriculation number: & 4756033 \\[0.6em]
Project title: & Confidence, not conflict: which value$\to$RT mapping best describes probability discounting, and does RT help? \\[0.6em]
Group members, if applicable: & Ashhad Raza Quadri, Saliq Neyaz, Maria Alejandra Pabon Galindo \\
\end{tabularx}

\bigskip

\subsection*{AI use statement}

Please tick one option:

\medskip

\noindent
$\square$ \textbf{No AI tools used.}

\noindent
I confirm that I did not personally use generative AI tools, AI-based writing
assistants, translation tools, image generators, code assistants, or similar AI
systems in preparing my contribution to this project.

\medskip

\noindent
$\blacksquare$ \textbf{AI tools used and disclosed below.}

\noindent
I confirm that I used AI tools only in limited and transparent ways that preserved
the independent character of my work. I remain fully responsible for the accuracy,
academic integrity, citations, methodology, code, analyses, and final submitted
content.

\bigskip

\subsection*{Details of AI use}

If no AI tools were used, write ``Not applicable'' in the first row.

\medskip

\begin{tabularx}{\textwidth}{|
>{\raggedright\arraybackslash}X|
>{\raggedright\arraybackslash}X|
>{\raggedright\arraybackslash}X|
>{\raggedright\arraybackslash}X|}
\hline
\textbf{AI tool used} &
\textbf{Purpose of use} &
\textbf{Parts of the project affected} &
\textbf{Extent of use / student responsibility} \\
\hline
Claude Code (Anthropic 5.1) &
Assistance with writing and debugging Python code: drafting implementations of
routines I had specified, and diagnosing errors (optimiser convergence failures,
numerical problems in the shifted log-normal log-likelihood). &
Inference code for my components (Sec.~2.3): the MAP objective, the L-BFGS-B
wrapper with multiple restarts, the MLE cross-check and the likelihood-curvature
diagnostic. Reliability analysis code (Sec.~3.2, Table~3): Pearson/Spearman/ICC
computation and the bootstrap CI routines. Associated plotting code for
Figs.~2 and 5. &
AI produced initial drafts of these routines from specifications I wrote; the
choice of model, estimator, priors, diagnostics and statistics was mine.
I read, corrected and tested all generated code, verified its output against
the model-free estimates (Sec.~2.1) and the MLE cross-check, and can explain
every line. No AI was used to interpret results or to decide what the analyses
show. \\
[2.5em]
\hline
Claude Code (Anthropic 5.1) &
Language and grammar correction and rephrasing of text I had already drafted
myself. &
My own paragraphs in Sec.~2.3 (inference), Sec.~3.1--3.2 (identifiability and
reliability) and the corresponding passages of Sec.~4. &
Wording and clarity only. Every claim, number, argument and interpretation is my
own; no passage was accepted without my review, and no substantive prose was
generated from scratch by AI. \\
[2.5em]
\hline
\end{tabularx}

\bigskip

\subsection*{Prompt log}

\noindent
Did you retain a prompt log or notes documenting the AI-assisted steps?

\medskip

\noindent
$\square$ Yes, attached or available upon request.

\noindent
$\square$ No AI tools were used.

\noindent
$\square$ No prompt log was retained. Explanation:
\fillme{explanation}

\bigskip

\subsection*{Individual contribution and responsibility}

\noindent
Briefly describe your own contribution to the project, including report writing,
coding, modelling, analyses, figures, or interpretation of results:

\medskip

\noindent
I owned \textbf{RT model M2} (decision uncertainty, $\mu = a - b\,|2P-1|$): its
specification, psychological motivation, fitting and evaluation, including the
result that it gives the best run~B RT predictive log-likelihood ($-1.492$ nats
per trial) and is the only model ranking that holds both in sample on run~A and
out of sample on run~B. I also built the \textbf{inference machinery} (MAP
objective, priors and their justification in Table~1, the L-BFGS-B optimiser with
multiple restarts, the MLE cross-check and the likelihood-curvature
identifiability diagnostic, both in Fig.~2) and the \textbf{reliability analysis}
of Sec.~3.2, Fig.~5 and Table~3 (Pearson/Spearman/ICC, bootstrap CIs, paired CIs
on reliability differences), which the report reads as generalisation across
decision contexts rather than classical test--retest. I drafted the report text
for Sec.~2.3, Secs.~3.1--3.2 and the corresponding passages of Sec.~4.

\medskip

\noindent
The work was divided equally in three. Shared components -- the Rachlin
hyperboloid valuation model and logistic choice rule, the shifted log-normal RT
distribution, the model-comparison framework, the figures and this report -- were
developed jointly, and I reviewed and accept responsibility for the parts I
contributed to.

\bigskip

\noindent
In the case of a group project, I confirm that I have reviewed the shared report,
code, analyses, and figures to which I contributed. I accept responsibility for my
own contribution and for any AI-assisted material included in shared components
that I reviewed, used, or submitted as part of the project.

\bigskip

\subsection*{Final confirmation}

\noindent
I confirm that I have not submitted AI-generated content as my own without review,
verification, and disclosure. I have checked all factual claims, references, data,
code, analyses, and interpretations myself. I have not entered sensitive,
personal, confidential, or copyright-protected material into AI systems in
violation of data protection, copyright, or course rules.

\bigskip
\bigskip

\noindent
\textbf{Name and signature:} Saliq Neyaz
\hfill
\textbf{Date:} 30 August 2026

\end{document}
